\cotefile{157-5}
\scanpage{117}
\mspage{104}
Il reste à prouver (6.1.12), et on procède
comme pour 3.3 (p.~54), en utilisant
le fait qu'une flèche $f$ dans $M(I)_c$
qui est dans $W(I)_{\deleted{\unsure{c}}}$, s'écrit \inserted{dans $M(I)_c$} sous la
forme $f=pi$, où $i\in TC(I)$, et
où $p$ est \deleted{\unsure{l'inv.}} inverse à gauche (rétraction)
d'une flèche $j\in TC(I)$. On aura
donc $u_!(f)=u_!(p)u_!(i)$, \deleted{\unsure{et}} $u_!(i)\in TC(I')$
par (6.\unsure{1}.10), et $u_!(p)$ est inverse à
gauche de $u_!(j)\deleted{\unsure{\in}} \in TC(I')$. On a
donc $u_!(i)$, $u_!(j)$, $u_!(p)\in W(I')$, d'où
$u_!(f)\in W(I')$, qed.

Prouvons enfin d). \deleted{\unsure{Il}} Compte
tenu de la définition de $C_I$, pour
établir $C(I)\subset C_I$, il \deleted{\unsure{est} \illegible{1}} \inserted{suffit de}
voir ceci

% slanted
\margin{\uline{NB} Ce résultat se généralise sous les \unsure{condit\textsuperscript{ns}|conditions} \unsure{inscrites|précisées} dans 6.1.5.}
\uline{Corollaire} 6.2. Soit $I$ \unsure{quelc.} dans \unsure{Cat},
soient $U\subset V$ deux fermés dans $I$
(i.e. des ouverts dans $I^\circ$), et \deleted{\unsure{considé}}
soit $\deleted{f}\unsure{u|f}\colon X_\cdot\longrightarrow Y_\cdot$ une flèche
dans $\deleted{\unsure{f}}C(I)$ (resp. dans $TC(I)$).
\deleted{Considérons la \unsure{fl}} Posons alors
% archive sheet number: 116

